%% kapa: exam=1 round=21 year=2010 period=2 form=A subject=appraisal_law subject_ko=부동산관계법규 seq=17 number=57
%% source: Q-Net 감정평가사 기출문제 (cst00309) artlSeq=5024181 file=extracted/1008/5024181/2010년 제21회 감정평가사 1차시험 기출문제/2교시.hwp
%% conversion: kordoc 4.24.51 -> tools/gen1.py (auto); review=proofread-gpt-sol; flags=none
\begin{question}{57}
부동산 가격공시 및 감정평가에 관한 법령상 다음 보기 중 옳지 않은 것만으로 모두 묶인 것은?
\begin{bogi}
ㄱ. 감정평가업자는 자기 소유의 토지에 대하여는 감정평가하여서는 아니된다. \par
ㄴ. 감정평가업자는 감정평가서의 원본을 그 교부일로부터 5년간 보존하여야 한다. \par
ㄷ. 감정평가업자가 의뢰인으로부터 받는 수수료의 요율은 감정평가협회가 정한다. \par
ㄹ. 미성년자는 감정평가사가 될 수 없다. \par
ㅁ. 감정평가업자는 토지등의 매매업을 직접 영위하여서는 아니된다. \par
\end{bogi}
\begin{choices}
  \item ㄱ, ㄴ
  \item ㄱ, ㅁ
  \item ㄴ, ㄷ
  \item ㄷ, ㄹ
  \item ㄹ, ㅁ
\end{choices}
\end{question}
\answer{3}
